\chapter{Cortes y Secciones}

\section{Representación mediante Cortes y Secciones}

Cuando una pieza posee una estructura interna compleja, el uso exclusivo
de vistas exteriores con líneas ocultas (trazo discontinuo) puede
resultar confuso y difícil de interpretar. Para resolver este problema,
se emplean \textbf{cortes (o secciones)}, donde se imagina que la pieza
es seccionada por un plano imaginario, permitiendo visualizar y acotar
claramente los detalles interiores.

\subsection{El Plano de Corte}

El plano de corte es una superficie imaginaria utilizada para
seccionar la pieza. En el dibujo, este plano se representa mediante una
línea de trazo y punto grueso en los extremos, con flechas que indican
la dirección de observación.

\subsection{Tipos de Cortes}

Los cortes más comunes incluyen:
\begin{itemize}
	\item \textbf{Corte total:} El plano de corte atraviesa totalmente
	      la pieza.
	\item \textbf{Media sección:} Se utiliza para piezas simétricas,
	      mostrando la mitad en corte y la mitad en vista exterior.
	\item \textbf{Corte parcial (rotura):} Se corta solo una pequeña
	      parte de la pieza para revelar un detalle específico.
	\item \textbf{Corte escalonado:} Se utiliza para piezas con
	      detalles internos alineados de forma compleja, usando
	      varios planos de corte paralelos.
\end{itemize}

\subsection{Normas de Representación: El Rayado}

Las superficies que han sido atravesadas por el plano de corte deben
indicarse mediante un \textbf{rayado}.
\begin{itemize}
	\item El rayado consiste en líneas finas inclinadas a 45 grados
	      respecto a los ejes principales de la pieza.
	\item La separación entre las líneas debe ser uniforme y depender
	      del tamaño de la pieza.
	\item Ciertos elementos, como nervios de refuerzo, tornillos,
	      pernos, ejes o chavetas, \textbf{no se rayan} cuando el
	      plano de corte pasa por su eje longitudinal, incluso si
	      están siendo cortados.
\end{itemize}

La correcta aplicación de los cortes es esencial para la fabricación,
ya que permite definir dimensiones internas que, de otra forma, serían
imposibles de acotar con precisión.

\section{Cuestionario}
\begin{enumerate}
	\item ¿Qué es un plano de corte y cómo se representa?
	\item Define qué es una sección en dibujo técnico.
	\item ¿Cómo se representan las superficies cortadas (rayado)?
	\item ¿Cómo se indica la dirección de observación del corte?
	\item ¿Qué partes normalizadas no se deben rayar (ej: pernos)?
	\item Diferencia técnica entre corte total y media sección.
	\item ¿Para qué sirven los cortes parciales o roturas?
	\item ¿Cómo se identifica la sección resultante en el plano?
\end{enumerate}

\section{Parte Práctica}
\begin{enumerate}
	\item Dibujar una pieza simple y realizar un corte total longitudinal.
	\item Realizar un corte transversal a un cilindro hueco.
	\item Aplicar rayado normalizado (distancia y ángulo correctos) a
	      una pieza cortada.
	\item Dibujar un corte escalonado para mostrar dos detalles internos
	      de diferente profundidad.
	\item Realizar un corte parcial en una pieza con una parte interna simple.
	\item Dibujar una sección girada en el mismo plano de la vista principal.
	\item Crear una vista de sección para una pieza con nervios de
	      refuerzo (aplicando la regla de no rayar el nervio).
	\item Aplicar diferentes tipos de rayado según el material.
	\item Dibujar una pieza simétrica usando media sección.
	\item Realizar un plano que incluya una sección de un eje con chavetero.
\end{enumerate}
